% TOP_PROBLEM: {"id": 1500, "title": "Odd Weird Numbers", "category_id": 1, "category": {"id": 1, "name": "number_theory", "display_name": "Number Theory", "description": "Properties of integers, prime numbers, Diophantine equations.", "slug": "number-theory", "order_index": 1, "created_at": "2026-07-31T15:26:25.670Z"}, "status": "open"}
% FIRST_POSTED: null
% ENTRY_KIND: statement_only
% Imported from: batch08.tex; UnsolvedMath ID 1500
\documentclass[11pt]{article}
\usepackage[margin=25mm]{geometry}
\usepackage{fontspec}
\setmainfont{Latin Modern Roman}
\usepackage{amsmath,amssymb,mathtools}
\usepackage{unicode-math}
\setmathfont{Latin Modern Math}
\usepackage{xurl,hyperref}
\hypersetup{colorlinks=true,urlcolor=blue}
\setcounter{secnumdepth}{0}
\setlength{\parindent}{0pt}
\setlength{\parskip}{5pt}
\setlength{\emergencystretch}{3em}
\newcommand{\sref}[2]{\hyperlink{source:#1:#2}{[S#2]}}
\newcommand{\cataloguescope}[1]{\paragraph{Recorded target.}#1}
\begin{document}
% UnsolvedMath ID 1500; source catalogue code NUM-006
\setcounter{section}{1499}
\section{Odd Weird Numbers}\label{1500}
{\small\sffamily UnsolvedMath ID 1500 · Number Theory}
\subsection{Definitions and mathematical statement}

\paragraph{Shared notation.}$\mathbb N=\{1,2,\ldots\}$, and $\mathbb Z,\mathbb Q,\mathbb R,\mathbb C$ denote the integers, rationals, reals and complex numbers. Any other conventions are specified locally.


For \(n\in\mathbb N\), let \(\sigma(n)\) be the sum of all
positive divisors of \(n\), and put
\(D^*(n)=\{d\in\mathbb N:d\mid n,\ d<n\}\).
The integer \(n\) is abundant if \(\sigma(n)>2n\).
It is semiperfect (also called pseudoperfect) if some subset
\(S\subseteq D^*(n)\) satisfies \(\sum_{d\in S}d=n\).
It is weird if it is abundant and not semiperfect.
Divisors in \(S\) occur at most once.
\paragraph{Statement recorded in the supplied source.}
Does there exist an odd integer \(n\ge1\) such that
\[
 \sigma(n)>2n\quad\text{and}\quad
 \sum_{d\in S}d\ne n\ \text{for every subset }S\subseteq D^*(n)?
\] \sref{1500}{2}
\subsection{Short English statement}
Can an odd integer have more than itself as the sum of its proper divisors, while no subset of those divisors sums exactly to it?
\subsection{Sources}
\begin{description}
\item[\hypertarget{source:1500:1}{S1}] ULAM.AI and the UnsolvedMath Contributors, \emph{UnsolvedMath: A Curated Collection of Open Mathematics Problems}, record 1500 (NUM-006). CC BY 4.0. \url{https://huggingface.co/datasets/ulamai/UnsolvedMath}.
\item[\hypertarget{source:1500:2}{S2}] Wenjie Fang, \emph{Searching on the boundary of abundance for odd weird numbers}, arXiv:2207.12906 (2022), Section 1, abundance, semiperfectness and the odd-weird-number existence question. \url{https://arxiv.org/abs/2207.12906}.
\end{description}
\label{end:1500}
\end{document}
